% ══════════════════════════════════════════════════════════════════════════
% Recap frames. Called as checkpoints between blocks, and
% the closing frame at the very end. Kept in one file so the wording of the
% recaps can be kept consistent.
% ══════════════════════════════════════════════════════════════════════════

% ── Checkpoint after the count cases ─────────────────────────────────────
\newcommand{\recapdch}{%
\begin{frame}{dCH So Far, in Four Lines}

  \begin{enumerate}
    \setlength{\itemsep}{7pt}
    \item \textbf{Switchers and stayers.} Each effect compares groups whose
      treatment changes at $F_g$ with groups that held the \emph{same}
      period-one treatment and had not yet moved.
    \item \textbf{Effect$_\ell$ and Placebo$_\ell$.} The event study reads
      forward from the switch; the placebos read backward and are zero under
      parallel trends and no anticipation.
    \item \textbf{Normalization.} With a non-binary treatment, \texttt{normalized}
      turns the effect into an effect \emph{per dose-period accumulated up to $\ell$}, and
      \texttt{normalized\_weights} shows how much of it comes from the current
      treatment and how much from its lags.
    \item \textbf{Two count cases.} Favara \& Imbs and Gentzkow et al.: the same
      command, clean placebos in both.
  \end{enumerate}

  \medskip
  Next: designs where the placebos \emph{fail}, treatments that arrive in
  sequence, and treatments that are a dose rather than a switch.

\end{frame}%
}

% ── Checkpoint before synthetic control ──────────────────────────────────
\newcommand{\recapbeforescm}{%
\begin{frame}{Where We Have Got To}

  \begin{columns}[T]
    \begin{column}{0.48\textwidth}
      \begin{block}{de Chaisemartin \& D'Haultf\oe{}uille}
        \begin{itemize}\small
          \item The control group is a set of \textbf{stayers with the same
            baseline treatment}.
          \item Identification comes from \textbf{many units switching at
            different dates}.
          \item Parallel trends is imposed on the \textbf{status-quo} outcome,
            and the placebos test it.
        \end{itemize}
      \end{block}
    \end{column}
    \begin{column}{0.48\textwidth}
      \begin{alertblock}{What changes now}
        \begin{itemize}\small
          \item \textbf{One} treated unit, or a handful.
          \item No cohorts, no group-time effects, no asymptotics in the
            number of treated units.
          \item The control group has to be \textbf{built}, and its credential
            is the fit before treatment.
        \end{itemize}
      \end{alertblock}
    \end{column}
  \end{columns}

  \medskip
  The question stays the same throughout: \textbf{who is the control, and what
  exactly is being compared?}

\end{frame}%
}

% ── Closing frame of the course ──────────────────────────────────────────
\newcommand{\closingframe}{%
\begin{frame}{Who Is the Control: Five Methods}

  \begin{center}
  \footnotesize
  \renewcommand{\arraystretch}{1.35}
  \begin{tabular}{@{}p{2.9cm}p{4.5cm}p{4.6cm}@{}}
    \toprule
    \textbf{Method} & \textbf{Who is the control} & \textbf{What you must defend} \\
    \midrule
    dCH & Stayers with the \emph{same} period-one treatment
        & Parallel trends on the status-quo outcome; no crossing \\
    dCH \texttt{continuous()} & No two groups share a baseline, so the
        status-quo trend is modelled as a polynomial in it
        & The order of that polynomial \\
    Callaway, Goodman-Bacon \& Sant'Anna & Units with $D=0$, dose by dose
        & Parallel trends for levels; \emph{strong} parallel trends before you
          read any slope \\
    Synthetic control & A convex combination of donors that reproduces the
        pre-period
        & Fit before treatment, and a donor pool that can produce it \\
    Synthetic DiD & The same, reweighted over units \emph{and} periods
        & The same, with inference that suits the number of treated units \\
    \bottomrule
  \end{tabular}
  \end{center}

  \medskip
  \footnotesize No estimator removes the need to say which comparison you are
  making. Each of these makes the comparison explicit, which is what lets a
  reader disagree with it.

\end{frame}%
}
